\section{Bitcoin and Python Compared}
\label{sec:compare}

The Python study designed the comparison this paper completes. Table~\ref{tab:compare} sets the two projects side by side on the same detector, using Python's extended corpus (1999--2026, 299,192 candidates) and Bitcoin's (2011--2026, 62,791 candidates).

\begin{table}[t]
\caption{Bitcoin and Python on the same instrument (\% of candidates unless stated).}
\label{tab:compare}
\small
\setlength{\tabcolsep}{4pt}
\begin{tabular}{lrr}
\toprule
 & Bitcoin & Python \\
\midrule
Proposals with candidates & 199 & 764 \\
Candidate sentences & 62,791 & 299,192 \\
Participants & 636 & 3,639 \\
Mechanisms per candidate & 1.38 & 1.27 \\
\midrule
Ecosystem & 33.7 & 18.4 \\
Operational & 22.0 & 27.0 \\
Security & 17.5 & 2.6 \\
Economic & 14.8 & 2.3 \\
Compatibility & 12.1 & 13.6 \\
Functional & 7.6 & 25.8 \\
Standards & 6.5 & 8.8 \\
Organizational & 4.5 & 3.2 \\
Strategic & 4.2 & 5.0 \\
Tactical & 3.3 & 4.9 \\
User demand & 3.1 & 3.6 \\
Authority & 2.5 & 7.6 \\
Coalition & 2.1 & 2.0 \\
\midrule
Supporting / blocking / revising & 1.5 / 1.4 / 8.7 & 2.0 / 1.5 / 8.8 \\
External or mixed scope & 27.6 & 7.7 \\
Formal leader's share of candidates$^a$ & 1.0 & 8.0 \\
Community members' share & 57.6 & 37.1 \\
Any controversial mechanism & %%CONTRO_PCT%% & 4.7 \\
Supporting aligned with outcome & 74.9 & 70.0 \\
Blocking aligned with outcome & 12.1 & 26.5 \\
\bottomrule
\multicolumn{3}{l}{\footnotesize $^a$ lead maintainer (Bitcoin); BDFL plus Steering Council (Python).}
\end{tabular}
\end{table}

\paragraph{What the taxonomy predicted.} The Python paper expected Bitcoin to show stronger ecosystem, economic, security and deployment-based influence and weaker internal role-based influence. All four expectations hold, and by large margins: ecosystem signals are 1.8 times as frequent, security 6.7 times, economic 6.3 times, and authority a third as frequent. Scope confirms the same picture from another angle: 27.6\% of Bitcoin's candidates name external constituencies---users, wallets, exchanges, miners, companies---against 7.7\% in Python. Bitcoin's arguments are about the world outside the repository because the world outside the repository decides.

\paragraph{What it did not predict.} Three differences were not anticipated. First, \emph{functional} argument---whether the design solves the problem---is Python's second mechanism (25.8\%) and Bitcoin's sixth (7.6\%). Python's discussions are about language design, where the question is whether a feature is worth having; Bitcoin's are about protocol changes, where the question of whether a change is desirable is settled quickly and the argument moves to whether it is safe, what it costs and who will run it. Second, the formal leader is nearly absent from Bitcoin's record. Python's BDFL wrote 8\% of candidates and pronounced on most contested PEPs; Bitcoin's lead maintainers wrote 1\% and, after 2015, nothing that the tool counts as argument. The lead maintainer role in Bitcoin was administrative---merge, release, announce---and its holders declined to argue proposals on the list, which is why the era comparison shows no analogue of Python's authority-invoking council. Third, the Python paper's temporal finding---signalling peaking in the four weeks around the decision---does not transfer, for a reason that is itself a finding about governance: in Bitcoin the repository's status change records a deployment, not a decision, and the decision is distributed across merges, releases, signalling periods and node upgrades that no single timestamp captures.

\paragraph{What is the same.} Direction distributions are nearly identical (about 88\% neutral, 9\% revising, 1.5\% supporting, 1.4\% blocking), stance alignment with outcomes sits at the base rate in both projects, coalition signals are equally rare (2\%), and controversial mechanisms mark about one candidate in twenty-five in Bitcoin and one in twenty in Python. Two communities with opposite governance designs---a single decision-maker replaced by a council, and no decision-maker at all---argue with the same explicitness and the same small share of contested moves; what differs is what they argue about and who does the arguing.

